\documentclass[11pt]{article}
\usepackage[spanish]{babel}
\usepackage{amsmath}  % Math
\usepackage{amssymb}  % Symbols
\usepackage{graphicx} % Images
\graphicspath{ {./images/} }
\usepackage[utf8]{inputenc}
\usepackage[T1]{fontenc}
\usepackage[margin=1in]{geometry}
\usepackage{fancyhdr} % Para el pie de página personalizado
\usepackage{hyperref} % Para metadatos y enlaces en el PDF

% ----- CONFIGURACIÓN DE AUTORÍA Y PIE DE PÁGINA -----
\hypersetup{
    pdftitle={Tu Título},
    pdfauthor={Felipe Colli Olea},
    colorlinks=true,
    linkcolor=black,
    urlcolor=cyan
}

% Configuración del pie de página con fancyhdr
\pagestyle{fancy}
\fancyhf{} % Limpia la cabecera y el pie de página actuales
\fancyfoot[C]{\footnotesize Felipe Colli, Todos los Derechos Reservados} % Tu texto en el centro
\fancyfoot[R]{\thepage} % Número de página a la derecha
\renewcommand{\headrulewidth}{0pt} % Elimina la línea de la cabecera

\title{Trabajo de Lenguaje}

\date{\today}

\begin{document}

\maketitle
\thispagestyle{fancy} % Aplica el estilo de pie de página a la primera página también

\section*{Texto 1: Pregunta Integradora}

La literatura universal, en sus diversas manifestaciones, ofrece un espejo en el que se reflejan las complejidades del ser humano y las sociedades que este construye. A través del análisis de \textit{"El asesinato"} de Antón Chéjov, \textit{"El árbol"} de María Luisa Bombal y \textit{"Funes, el memorioso"} de Jorge Luis Borges, es posible desentrañar visiones del mundo que, aunque distantes en geografía y estilo, convergen en la exploración de personajes atrapados por sus circunstancias, ya sean sociales, emocionales o existenciales.

En \textbf{\textit{"El asesinato"}}, Chéjov nos sumerge en la cruda realidad de la Rusia zarista para exponer una sociedad que deshumaniza a través de la explotación. La visión del ser humano es la de un individuo reducido a una simple fuerza de trabajo. La joven Varka, privada del sueño y sometida a una labor incesante, pierde su capacidad de empatía y su juicio. El llanto del bebé se transforma en la manifestación de su opresión, un "enemigo" que le impide el descanso vital. El conflicto es, por tanto, eminentemente social: la indiferencia de sus amos y la precariedad de su existencia la empujan a un colapso psicológico que culmina en una tragedia. La sociedad es retratada como una máquina indiferente que consume la inocencia y la vida.

Por su parte, María Luisa Bombal en \textbf{\textit{"El árbol"}} traslada el conflicto al universo íntimo y psicológico de su protagonista. Brígida representa a la mujer relegada a un rol ornamental dentro de una sociedad patriarcal. Su conflicto no es la explotación física, sino el vacío emocional de un matrimonio sin amor. La visión del ser humano es la de un ser en busca de identidad y conexión, que se refugia en un mundo interior para sobrevivir. El gomero se convierte en el símbolo de esa vida interior, un santuario que la aísla de la mediocridad de su realidad. Cuando el árbol es derribado, la ilusión se desvanece y Brígida se ve forzada a confrontar la verdad de su vida, encontrando en la ruptura su única vía de liberación.

Finalmente, Jorge Luis Borges en \textbf{\textit{"Funes, el memorioso"}} nos presenta un conflicto puramente existencial y filosófico. La visión del ser humano se explora a través de los límites de la percepción y el pensamiento. Ireneo Funes, con su memoria infalible, es un ser anómalo que experimenta la realidad de una forma abrumadora. Es incapaz de olvidar, y por lo tanto, incapaz de pensar en términos abstractos o generales. Su mente es un archivo infinito de detalles inconexos, lo que le impide comprender el mundo de una manera conceptual. A diferencia de los otros cuentos, la sociedad es un mero telón de fondo; el verdadero conflicto ocurre dentro de la mente de Funes, cuya condición lo aísla radicalmente de una experiencia humana "normal".

\subsection*{Comparación, Similitudes y Diferencias}

Al comparar estos tres relatos, encontramos una similitud fundamental: \textbf{todos sus protagonistas son prisioneros}. Varka es prisionera de la explotación social; Brígida es prisionera de un matrimonio asfixiante y de las convenciones sociales; y Funes es prisionero de su propia mente prodigiosa. En los tres casos, esta condición de encierro conduce a una forma de parálisis o a un desenlace trágico: el acto desesperado de Varka, la postergada liberación de Brígida y la parálisis intelectual de Funes.

Sin embargo, las diferencias son cruciales y definen el enfoque de cada autor. La principal divergencia radica en la \textbf{naturaleza del conflicto y la representación de la realidad}.
\begin{itemize}
    \item \textbf{Chéjov} presenta un conflicto \textbf{externo y social} a través de un \textbf{realismo crudo y directo}. Su prosa es funcional y busca denunciar una injusticia material y visible.
    \item \textbf{Bombal} se centra en un conflicto \textbf{interno y emocional}, utilizando una \textbf{prosa lírica y subjetiva}. La realidad se filtra a través de la sensibilidad de Brígida, cargada de símbolos y atmósferas oníricas.
    \item \textbf{Borges}, en cambio, aborda un conflicto \textbf{intelectual y filosófico}. Su estilo es preciso, casi ensayístico, y utiliza un elemento fantástico (la memoria perfecta) no para evadir la realidad, sino para interrogarla en un plano abstracto.
\end{itemize}
Así, mientras el cuento ruso pone el foco en la opresión del cuerpo y la clase social, los relatos latinoamericanos exploran opresiones más íntimas: la del espíritu femenino en Bombal y la de la conciencia en Borges. Juntos, ofrecen un panorama de cómo las jaulas que limitan al ser humano pueden ser tanto materiales y sociales como psicológicas y existenciales.

\vspace{1cm}

\section*{Texto 2: Crítica Social y Poder}

La literatura latinoamericana ha explorado con particular agudeza las complejas dinámicas de poder que estructuran sus sociedades. A través de los cuentos \textit{"Muerte constante más allá del amor"} de Gabriel García Márquez, \textit{"Corrido"} de Juan José Arreola y \textit{"Ana María"} de José Donoso, se revelan distintos mecanismos de dominación —políticos, culturales y sociales— y su profundo impacto en la vida de los personajes, demostrando cómo el poder, en sus múltiples formas, puede corromper, aislar y deshumanizar.

En \textbf{\textit{"Muerte constante más allá del amor"}}, Gabriel García Márquez disecciona el \textbf{poder político} en su faceta más demagógica y vacía. El senador Onésimo Sánchez ejerce su dominio a través de un circo de promesas ilusorias, construyendo fachadas de cartón para ocultar la miseria de su pueblo. Este mecanismo de poder se basa en la manipulación y el engaño para perpetuar un sistema corrupto del que él mismo se beneficia. Sin embargo, el relato muestra la ironía de su condición: a pesar de su poder público, el senador es un hombre profundamente solo e impotente ante su propia muerte. Su poder no puede comprarle ni tiempo ni afecto genuino, lo que lo lleva a un acto final de vulnerabilidad y desesperación. El cuento es una crítica a la soledad del poder y a la brecha insalvable entre el discurso político y la realidad humana.

Por su parte, Juan José Arreola en \textbf{\textit{"Corrido"}} expone el \textbf{poder cultural} del machismo y los rígidos códigos de honor. El mecanismo de dominación no es político, sino social y simbólico. Los dos rivales no luchan por amor, sino por la afirmación de su masculinidad frente a la comunidad. La mujer es despojada de su agencia y convertida en un mero objeto, el pretexto para un duelo a muerte que es, en realidad, un ritual de validación masculina. Los personajes son víctimas de este sistema de poder cultural que los conduce fatalmente a la violencia. El efecto es devastador: los hombres mueren para cumplir con el guion social, y la mujer, aunque sobrevive, queda marcada por la "mala fama", estigmatizada por un conflicto que ella no inició. Arreola critica así una estructura cultural que sacrifica vidas en el altar del honor viril.

Finalmente, José Donoso en \textbf{\textit{"Ana María"}} aborda una forma de poder más sutil pero igualmente destructiva: el de la \textbf{negligencia y la marginación social}. El poder aquí se manifiesta en la indiferencia. Los padres de Ana María ejercen su poder sobre ella a través del abandono, tratándola como un objeto sin importancia. A su vez, la sociedad ejerce su poder sobre el anciano protagonista al marginarlo por su vejez e inutilidad laboral. El efecto de esta dinámica es la creación de un vínculo inesperado entre los dos seres más despojados de poder. Su relación se convierte en un acto de resistencia contra la indiferencia del mundo. Donoso critica una sociedad que ha roto sus lazos más básicos de cuidado y responsabilidad, donde el afecto solo parece posible en los márgenes, entre aquellos que han sido olvidados y despojados de todo poder.

En conclusión, estos tres relatos ofrecen una visión panorámica y crítica de las estructuras de poder en Latinoamérica. García Márquez expone la farsa del poder político, Arreola denuncia la tiranía del poder cultural machista, y Donoso revela la crueldad del poder que emana de la indiferencia social. En todos ellos, el poder se muestra como una fuerza que distorsiona las relaciones humanas, genera violencia y soledad, y condena a sus personajes a destinos marcados por la tragedia y la marginación.\end{document}
